% Part: lambda-calculus
% Chapter: introduction
% Section: fixed-point-combinator

\documentclass[../../../include/open-logic-section]{subfiles}

\begin{document}

\olfileid{lam}{int}{fix}
\olsection{固定点组合子}

假设有一个λ-项 $g$，且希望找到另一个项 $k$，使得 $k$ 与 $gk$ β-等价。
按我们的记法约定，定义项
\[
\fn{diag}(x) = xx
\]
与
\[
l(x) = g(\fn{diag}(x))
\]
换言之，$l$ 即为项
$\lambd[x][g(xx)]$。令 $k$ 为项 $ll$。则有
\begin{align*}
k & = (\lambd[x][g(xx)])(\lambd[x][g(xx)]) \\
& \red  g((\lambd[x][g(xx)])(\lambd[x][g(xx)])) \\
& = gk.
\end{align*}
若取
\[
Y = \lambd[g][((\lambd[x][g(xx)])(\lambd[x][g(xx)]))]
\]
则 $Yg$ 和 $g(Yg)$ 归约至同一项；故 $Yg \equiv_\beta g(Yg)$。
这被称为“柯里组合子”。反之，若取
\[
Y = (\lambd[xg][g(xxg)])(\lambd[xg][g(xxg)])
\]
则事实上 $Yg$ 归约至 $g(Yg)$，这是一个更强的断言。
$Y$ 的后一种形式被称为“图灵 组合子”。

\end{document}